% Generated from validated synthesis. Do not hand-edit.
\section*{Monthly Signals}
\addcontentsline{toc}{section}{Monthly Signals}
\smallskip
\noindent\textbf{Frontier Models \& Voice} model更新はdefault experience、API lifecycle、voiceという利用surfaceの更新と一体化した。 \hfill{\footnotesize p.~\pageref{pkg:frontier-models-may}}\par
\smallskip
\noindent\textbf{Agents \& Runtime} agentを長く動かすため、sandbox、approval、network policy、CLI／MCP統合が実行環境の主要設計項目になった。 \hfill{\footnotesize p.~\pageref{pkg:agents-runtime-may}}\par
\smallskip
\noindent\textbf{Inference \& Serving} KV cache、MoE、kernel、runtimeの各層が同時に更新され、推論最適化がstack全体の問題として見えるようになった。 \hfill{\footnotesize p.~\pageref{pkg:inference-serving-may}}\par
\smallskip
\noindent\textbf{Agent Safety \& Security} 攻撃面はpromptだけでなくworkspace、long-term memory、retrieval、monitoring policyへ広がった。 \hfill{\footnotesize p.~\pageref{pkg:agent-safety-security-may}}\par
\smallskip
\noindent\textbf{Capability \& Evaluation} benchmark外の研究成果とthird-party evaluation設計が、frontier capabilityをどう確かめるかという問題を前景化した。 \hfill{\footnotesize p.~\pageref{pkg:capability-evaluation-may} / p.~\pageref{pkg:paper-watch-may}}\par
\medskip
\begin{claimboundary}[Retrospective scope]
本号は2026年7月を後日確認可能になった一次情報も用いて再構成するRetrospective Specialである。Coverage Periodと制作時点を同一視せず、vendor / project / author claimの境界は各記事とTechnical Notesで明示する。
\end{claimboundary}
\medskip
\tableofcontents
